\documentclass[letterpaper,11pt]{article}

\usepackage[top=0.55in, bottom=0.55in, left=0.65in, right=0.65in]{geometry}
\usepackage{enumitem}
\usepackage{titlesec}
\usepackage[hidelinks]{hyperref}
\usepackage[T1]{fontenc}
\usepackage{lmodern}
\usepackage{microtype}
\usepackage{needspace}

\pagestyle{empty}
\setlength{\parindent}{0pt}
\setlist[itemize]{
  itemsep=1pt,
  topsep=3pt,
  parsep=0pt,
  leftmargin=1.25em,
  label=\textbullet
}

\titleformat{\section}
  {\bfseries\normalsize}
  {}
  {0em}
  {}
  [\vspace{1pt}\titlerule]
\titlespacing*{\section}{0pt}{9pt}{4pt}

% #1 org  #2 location/tech  #3 role/subtitle  #4 dates
\newcommand{\entry}[4]{%
  \par\vspace{5pt}
  \noindent\textbf{#1}\hfill\textit{#2}%
  \ifx\relax#3\relax\else%
    \\[-1pt]\textit{#3}\hfill\textit{#4}%
  \fi%
  \par\vspace{1pt}%
}

\begin{document}

%------------------------------------------------------------------
%  HEADER
%------------------------------------------------------------------
\begin{center}
  {\LARGE\bfseries Rithul Kamesh}\\[5pt]
  \small
  \href{mailto:hi@rithul.dev}{hi@rithul.dev}\enspace\textbar\enspace
  \href{https://rithul.dev}{rithul.dev}\enspace\textbar\enspace
  \href{https://github.com/rithulkamesh}{github.com/rithulkamesh}\enspace\textbar\enspace
  +91\,8904973890
\end{center}

\vspace{2pt}

%------------------------------------------------------------------
\section{SUMMARY}

\noindent
AI infrastructure and systems engineer building the runtime, storage, and tooling
layer beneath LLM and ML workloads. Shipped a multi-tier LLM caching runtime
achieving 92.5\% token reduction on agent workloads, and an embedded vector
database distributed across five language ecosystems, both running in production
AI pipelines. Systems background in Rust, Go, C++, and CUDA spanning execution
engines, distributed storage, and developer tooling.

%------------------------------------------------------------------
\section{SELECTED PROJECTS}

\entry{Continuum}{C++, Python}
{Unified LLM and ML Execution Runtime}{}
\begin{itemize}
  \item Designed a unified IR dispatching LLM token generation and tensor computation
        through a single C++ execution graph, with a five-tier reuse hierarchy --- exact
        memoization, semantic caching, trie prefix KV caching, layer KV warm-start, and
        memory graph recall --- achieving 92.5\% token reduction on mixed agent workloads
        and up to 99\% on shared-prefix calls
  \item Built provider-agnostic backend support (OpenAI, Azure OpenAI, Anthropic, vLLM,
        libtorch, Apple MLX) with durable, checkpointable execution enabling
        process-resumable and forkable agent workflows
  \item Shipped as \texttt{continuum-ai} on PyPI with Python/C++ bindings, CI matrix
        on Linux and macOS, fuzz workflow, and coverage gates; featured on Product Hunt
        at launch
\end{itemize}

\entry{Veclite}{Rust, Python (PyO3), Go (CGO), C/C++, Zig}
{Embedded Vector Database}{}
\begin{itemize}
  \item Engineered HNSW indexing with SIMD-accelerated distance kernels (cosine, L2,
        dot product, Manhattan) in Rust with single-file storage and a vector-native SQL
        dialect supporting hybrid metadata + similarity search
  \item Maintained 14 stable releases (latest v1.1.1) with an ABI-stable C layer,
        distributed across 5 language ecosystems --- Crates.io, PyPI (GIL-released
        search via PyO3), Go modules, and C/C++ headers --- enabling zero-server,
        in-process vector search for production AI memory pipelines
\end{itemize}

\entry{Docproc}{Python}
{Document Intelligence Engine}{}
\begin{itemize}
  \item Built a document-to-Markdown conversion engine preserving layout, tables, and
        semantic structure for LLM ingestion pipelines, with vision-model-based figure
        and diagram description
  \item Shipped an evaluation harness with a benchmarking suite measuring downstream
        reasoning fidelity across document types
\end{itemize}

%------------------------------------------------------------------
\section{EXPERIENCE}

\entry{H2LooP}{On-Site}
{AI Developer Advocate Intern}{July 2026 -- October 2026}
\begin{itemize}
  \item Built technical demos and reference implementations on Hydron's AI-powered
        embedded development platform, covering firmware generation, debugging,
        cross-target porting, and hardware-aware code generation workflows
  \item Produced developer-facing technical content and sample repositories, and ran
        hands-on workshops introducing AI-assisted embedded workflows to university and
        professional engineering audiences
\end{itemize}

%------------------------------------------------------------------
\Needspace{7\baselineskip}
\section{TECHNICAL SKILLS}

\noindent
\textbf{Languages:} Rust, Go, Python, C/C++, CUDA \\
\textbf{Systems:} Distributed Systems, SIMD, Lock-free Concurrency, WGPU, OpenMP, FDTD \\
\textbf{Infrastructure:} Docker, Kubernetes, Prometheus, Grafana, OpenTelemetry, CI/CD \\
\textbf{Backend:} REST, gRPC, FastAPI, Redis, Redis Streams, PostgreSQL, pgvector,
  BM25, WebSockets \\
\textbf{AI / ML:} PyTorch, Transformers, LangChain, LangGraph, Agent Systems,
  RAG, LLM Evaluation, HNSW, Vector Databases

%------------------------------------------------------------------
\section{EDUCATION}

\entry{PES University}{Bengaluru, India}
{B.Tech, Electronics and Communication Engineering}{Aug 2023 -- May 2027}

\entry{BGS National Public School}{Bengaluru, India}
{HSSC, SSC}{Feb 2023}

%------------------------------------------------------------------
\section{RESEARCH EXPERIENCE}

\entry{Photonics and Quantum Technology Lab, PES University}{Bengaluru, India}
{Undergraduate Research Assistant}{Feb 2025 -- Present}
\begin{itemize}
  \item Modeling photonic integrated circuits via Lindblad/Hamiltonian open-quantum-system
        formalisms and building FDTD-coupled quantum solvers (Crank--Nicholson stepping,
        CPML boundaries, NTFF transforms) for photonic device characterization and state
        fidelity analysis
\end{itemize}

\vspace{3pt}
\noindent\textit{Teaching Assistant, Non-Linear Optics and Quantum Technology}
\hfill \textit{Jan 2026 -- May 2026}\\
Co-delivered lectures and built study guides and question banks for 100+ students.

%------------------------------------------------------------------
\section{PUBLICATIONS}

\noindent
B.\,S.\,Ullal, V.\,K.\,Narasimhaiah, and \textbf{R.\,Kamesh},
``ResNet based deep learning approach for chronic obstructive pulmonary disease
prediction using lung sound analysis,''
\textit{IAES International Journal of Artificial Intelligence (IJ-AI)} (Scopus),
vol.\,15, no.\,2, pp.\,1733--1745, Apr.\,2026.
[ResNet: 97.90\% accuracy, 98.72\% precision, 0.9937\,AUC]
\href{https://doi.org/10.11591/ijai.v15.i2.pp1733-1745}{doi:10.11591/ijai.v15.i2.pp1733-1745}

%------------------------------------------------------------------
\section{ACHIEVEMENTS}

\noindent
Hack Night 7.0, Third Place \hfill Oct 2025 \\
Speaker, Snowflake User Meetup Chennai \hfill Jun 2025 \\
Hack Night 6.0, Second Place \hfill Oct 2024

\end{document}
